\section{Определение предела функции в точке области. Связь с
существованием пределов действительной и мнимой части функции.
Определение непрерывной функции в точке и в
области.}\label{ux43eux43fux440ux435ux434ux435ux43bux435ux43dux438ux435-ux43fux440ux435ux434ux435ux43bux430-ux444ux443ux43dux43aux446ux438ux438-ux432-ux442ux43eux447ux43aux435-ux43eux431ux43bux430ux441ux442ux438.-ux441ux432ux44fux437ux44c-ux441-ux441ux443ux449ux435ux441ux442ux432ux43eux432ux430ux43dux438ux435ux43c-ux43fux440ux435ux434ux435ux43bux43eux432-ux434ux435ux439ux441ux442ux432ux438ux442ux435ux43bux44cux43dux43eux439-ux438-ux43cux43dux438ux43cux43eux439-ux447ux430ux441ux442ux438-ux444ux443ux43dux43aux446ux438ux438.-ux43eux43fux440ux435ux434ux435ux43bux435ux43dux438ux435-ux43dux435ux43fux440ux435ux440ux44bux432ux43dux43eux439-ux444ux443ux43dux43aux446ux438ux438-ux432-ux442ux43eux447ux43aux435-ux438-ux432-ux43eux431ux43bux430ux441ux442ux438.}

Пусть функция \(f(z)\) определена в проколотой окрестности точки
\(z_0\), т.е. в кольце \[K: 0 < |z - z_0| < R.\] \textbf{Определение
предела по Коши:} Комплексное число \(A\) называется пределом функции
\(f(z)\) при \(z \to z_0\), если
\[\forall \varepsilon > 0\ \exists \delta > 0: \forall z,\ 0 < |z - z_0| < \delta \Rightarrow |f(z) - A| < \varepsilon.\]
Обозначение: \(\lim\limits_{z \to z_0} f(z) = A\), или \(f(z) \to A\)
при \(z \to z_0\).

\textbf{Определение предела по Гейне:} Комплексное число \(A\)
называется пределом функции \(f(z)\) в точке \(z_0\), если для любой
последовательности \(\{z_n\}\), \(z_n \in K\), сходящейся к \(z_0\),
последовательность \(\{f(z_n)\}\) сходится к \(A\):
\[\forall \{z_n\},\ z_n \in K,\ \lim_{n\to\infty} z_n = z_0 \;\Rightarrow\; \lim_{n\to\infty} f(z_n) = A.\]
Определения по Коши и по Гейне эквивалентны.

\textbf{Теорема (связь с Re/Im):} Пусть \(f(z) = u(x,y) + i\,v(x,y)\),
\(z_0 = x_0 + i y_0\). Существование предела
\(\lim\limits_{z \to z_0} f(z)\) равносильно существованию двух пределов
\(\lim\limits_{\substack{x \to x_0 \\ y \to y_0}} u(x,y)\) и
\(\lim\limits_{\substack{x \to x_0 \\ y \to y_0}} v(x,y)\), причём
\[\lim_{z \to z_0} f(z) = \lim_{\substack{x \to x_0 \\ y \to y_0}} u(x,y) + i \lim_{\substack{x \to x_0 \\ y \to y_0}} v(x,y).\]

\textbf{Определение:} \emph{Непрерывность функции:}

\begin{enumerate}
\def\labelenumi{\arabic{enumi}.}
\tightlist
\item
  в точке
  \(z=z_{0}\Leftrightarrow \lim\limits_{ z \to z_{0} }f(z)=f(z_{0})\Leftrightarrow\Delta f=f(z_{0}+\Delta z)-f(z_{0})=o(1)\Leftrightarrow \lim\limits_{ \Delta z \to 0 }\Delta f=0\)
\item
  В области \(G\) \(\Leftrightarrow f(z)\) непрерывна в
  \(\forall z_{0}\in G\)
\item
  в замкнутой области \(\overline{G}\) - непрерывная \(\forall z_{0}\)
  внутренних + \(\forall z_{0}\in\partial G\)
\end{enumerate}

\textbf{Утверждение 1:} \(f(z)=u(x,y)+i\cdot v(x,y), \ z=x+iy\)
Непрерывность функции \(f(z)\) по определению выше \(\Leftrightarrow\)
непрерывность функций \(u(x,y)\) и \(v(x,y)\) по пунктам из определения
непрерывности.
